\section{代数的导子}

记号约定:$A$是分次$K$-代数,$\epsilon$是交换因子.

\begin{definition}{$\epsilon$-伴随}{}
	对$A$中的每个齐次元$a$,称$\ad_{\epsilon}\left(a\right):A\to A,x\mapsto\left[a,x\right]_{\epsilon}$是$a$的$\epsilon$-伴随.
\end{definition}

\begin{remark}{}{}
	$\ad_{\epsilon}\left(a\right)$是$\deg\left(a\right)$次分次$K$-线性映射.
\end{remark}

\begin{proposition}{$\epsilon$-伴随和导子的相容性}{}
	\begin{enumerate}
		\item 对每个$\epsilon$-导子$d:A\to A$和$A$的每个齐次元$a$,都有$\left[d,\ad_{\epsilon}\left(a\right)\right]_{\epsilon}=\ad_{\epsilon}\left(d\left(a\right)\right)$.
		\item 若$A$是结合的,则对$A$的每个齐次元$a$,它的$\epsilon$-伴随是$A$上的$\deg\left(a\right)$次$\epsilon$-导子.
	\end{enumerate}
\end{proposition}

\begin{definition}{$\epsilon$-内导子}{}
	设$A$是结合的.对$A$中的每个齐次元$a$,称$\ad_{\epsilon}\left(a\right)$是$a$的$\epsilon$-内导子.
\end{definition}

\begin{corollary}{内导子的李代数结构和Jacobi恒等式}{}
	设$A$是结合的.
	\begin{enumerate}
		\item 对$A$的任意齐次元$a,b$,有$\left[\ad_{\epsilon}\left(a\right),\ad_{\epsilon}\left(b\right)\right]_{\epsilon}=\ad_{\epsilon}\left(\left[a,b\right]_{\epsilon}\right)$.
		\item 设$a,b$是$A$的齐次元,次数分别为$\alpha,\beta$,则对每个$\gamma$和$A$的每个$\gamma$次齐次元$c$,Jacobi恒等式成立:
		\begin{align*}
			\epsilon\left(\alpha,\gamma\right)\left[a,\left[b,c\right]_{\epsilon}\right]_{\epsilon}+\epsilon\left(\beta,\alpha\right)\left[b,\left[c,a\right]_{\epsilon}\right]_{\epsilon}+\epsilon\left(\gamma,\beta\right)\left[c,\left[a,b\right]_{\epsilon}\right]_{\epsilon}=0.
		\end{align*}
	\end{enumerate}
\end{corollary}


























































































































































